% ============================================================
% Formulario, argomento 15: la gaussiana e la funzione Q (dispensa, capitolo 15)
% ============================================================
\argomento{La variabile gaussiana e la funzione $Q$}{Dispensa: capitolo 15}[Gaussiana e funzione Q]

\voce{gaussstd}{Gaussiana standard $\mathcal{N}(0, 1)$}{la campana centrata in 0}
  {f_{X_0}(x) = \frac{1}{\sqrt{2\pi}}\, e^{-\frac{x^2}{2}}}
  {\variante{media e varianza}{E[X_0] = 0, \quad \sigma_{X_0}^2 = E[X_0^2] = 1}
   \variante{massimo}{f_{X_0}(0) = \frac{1}{\sqrt{2\pi}} \simeq 0{,}399}}

\voce{gaussintegrale}{Integrale di Gauss}{in coordinate polari}
  {\int_{-\infty}^{\infty} e^{-\frac{x^2}{2}} \dd x = \sqrt{2\pi}}
  {}

\voce{gaussgen}{Gaussiana generica $\mathcal{N}(\mu_X, \sigma_X^2)$}{determinata da media e varianza}
  {f_X(x) = \frac{1}{\sqrt{2\pi\sigma_X^2}}\, e^{-\frac{(x - \mu_X)^2}{2\sigma_X^2}}}
  {\variante{dalla standard}{X = \sigma_X X_0 + \mu_X, \quad E[X] = \mu_X, \quad \sigma_X^2 = \mathrm{VAR}(X)}
   \variante{standardizzazione}{X_0 = \frac{X - \mu_X}{\sigma_X} \sim \mathcal{N}(0, 1)}}

\voce{funzioneq}{Funzione $Q$}{la coda della gaussiana standard}
  {Q(x) = P(X_0 > x) = \frac{1}{\sqrt{2\pi}} \int_x^{\infty} e^{-\frac{t^2}{2}} \dd t = 1 - F_{X_0}(x)}
  {\variante{la CDF della standard}{\Phi(x) = F_{X_0}(x) = 1 - Q(x)}}

\voce{codegauss}{Probabilità di una gaussiana qualsiasi}{$X \sim \mathcal{N}(\mu_X, \sigma_X^2)$}
  {P(X > x) = Q\left( \frac{x - \mu_X}{\sigma_X} \right)}
  {\variante{sotto una soglia}{P(X \leq x) = 1 - Q\left( \frac{x - \mu_X}{\sigma_X} \right)}
   \variante{intervallo}{P(a < X \leq b) = Q\left( \frac{a - \mu_X}{\sigma_X} \right) - Q\left( \frac{b - \mu_X}{\sigma_X} \right)}}

\voce{qproprieta}{Proprietà della funzione $Q$}{}
  {Q(-\infty) = 1 \qquad Q(0) = \frac{1}{2} \qquad Q(+\infty) = 0 \qquad Q(-x) = 1 - Q(x)}
  {\variante{è decrescente}{\frac{\dd Q(x)}{\dd x} = -\frac{1}{\sqrt{2\pi}}\, e^{-\frac{x^2}{2}} < 0}}

\voce{tavolaq}{Tavola della funzione $Q$}{valori calcolati, tre cifre \oltre}
  {\text{\footnotesize\setlength{\tabcolsep}{4pt}\renewcommand{\arraystretch}{1.15}%
   \begin{tabular}{@{}cl|cl|cl|cl@{}}
      $x$ & $Q(x)$ & $x$ & $Q(x)$ & $x$ & $Q(x)$ & $x$ & $Q(x)$ \\ \hline
      $0{,}0$ & $0{,}500$ & $1{,}0$ & $0{,}159$ & $2{,}0$ & $0{,}0228$ & $3{,}0$ & $0{,}00135$ \\
      $0{,}1$ & $0{,}460$ & $1{,}1$ & $0{,}136$ & $2{,}1$ & $0{,}0179$ & $3{,}1$ & $0{,}000968$ \\
      $0{,}2$ & $0{,}421$ & $1{,}2$ & $0{,}115$ & $2{,}2$ & $0{,}0139$ & $3{,}2$ & $0{,}000687$ \\
      $0{,}3$ & $0{,}382$ & $1{,}3$ & $0{,}0968$ & $2{,}3$ & $0{,}0107$ & $3{,}3$ & $0{,}000483$ \\
      $0{,}4$ & $0{,}345$ & $1{,}4$ & $0{,}0808$ & $2{,}4$ & $0{,}00820$ & $3{,}4$ & $0{,}000337$ \\
      $0{,}5$ & $0{,}309$ & $1{,}5$ & $0{,}0668$ & $2{,}5$ & $0{,}00621$ & $3{,}5$ & $0{,}000233$ \\
      $0{,}6$ & $0{,}274$ & $1{,}6$ & $0{,}0548$ & $2{,}6$ & $0{,}00466$ & $3{,}6$ & $0{,}000159$ \\
      $0{,}7$ & $0{,}242$ & $1{,}7$ & $0{,}0446$ & $2{,}7$ & $0{,}00347$ & $3{,}7$ & $0{,}000108$ \\
      $0{,}8$ & $0{,}212$ & $1{,}8$ & $0{,}0359$ & $2{,}8$ & $0{,}00256$ & $3{,}8$ & $7{,}23 \cdot 10^{-5}$ \\
      $0{,}9$ & $0{,}184$ & $1{,}9$ & $0{,}0287$ & $2{,}9$ & $0{,}00187$ & $3{,}9$ & $4{,}81 \cdot 10^{-5}$ \\
   \end{tabular}}}
  {\variante{per $x$ più grandi}{Q(4) \simeq 3{,}17 \cdot 10^{-5}, \quad Q(5) \simeq 2{,}87 \cdot 10^{-7}, \quad
                     Q(6) \simeq 9{,}87 \cdot 10^{-10}}}

\voce{qmaggiorazione}{Maggiorazione della funzione $Q$}{$x > 0$; per $x$ grande è anche un'approssimazione}
  {Q(x) < \frac{1}{x\sqrt{2\pi}}\, e^{-\frac{x^2}{2}}}
  {\variante{più semplice, per $x > \frac{1}{\sqrt{2\pi}} \simeq 0{,}4$}{Q(x) < e^{-\frac{x^2}{2}}}}

\voce{tresigma}{Regola delle tre sigma}{68--95--99,7}
  {P\left( \abs{X - \mu_X} \leq k\, \sigma_X \right) = 1 - 2\, Q(k)}
  {\variante{la coda doppia}{P\left( \abs{X - \mu_X} > k\, \sigma_X \right) = 2\, Q(k)}
   \variante{$k = 1, 2, 3$}{0{,}683, \quad 0{,}954, \quad 0{,}997}}

\voce{erroregauss}{Errore di una trasmissione binaria}{$\pm A$ equiprobabili, rumore $\mathcal{N}(0, \sigma^2)$, soglia 0}
  {P_e = Q\left( \frac{A}{\sigma} \right)}
  {}
